\section{Introduction}
In recent years, skills have emerged as an increasingly popular paradigm for extending LLM agents. This paradigm was originally introduced by Anthropic~\cite{agentskills-overview} as an open skill-packaging convention and has since been adopted by a growing number of agent frameworks, such as OpenAI~\cite{openai-codex-skills}, LangChain~\cite{langchain-deepagents-skills}, and OpenClaw~\cite{openclaw-skills-docs}. Specifically, instead of embedding all capabilities into a single agent, skills bundle prompts, code scripts, and configurations into reusable modules that can be installed and composed on demand.
As this paradigm has gained traction, a growing number of public registries and marketplaces have emerged to support skill publication, discovery, installation, and reuse. Representative examples include OpenClaw’s ClawHub~\cite{clawhub}, \texttt{skills.sh}~\cite{skills.sh}, and \texttt{SkillsMP}~\cite{skillsmp}. As of March, 2026, ClawHub had hosted over 30K skills~\cite{clawhub}, while \texttt{SkillsMP} had collected over 546K~\cite{skillsmp}. These developments indicate that skills are rapidly evolving into a new form of \emph{agentic supply chain}~\cite{hu2025llmsc,wang2025llmsc,huang2025llmsc}, where third-party skills are created, distributed, discovered, installed, and reused in the agentic ecosystems.

However, the widespread adoption of skills also introduces a new attack surface into the emerging agentic supply chain. In platforms such as OpenClaw~\cite{openclaw-skills-docs}, skills are used to teach the agent how to invoke tools, while the underlying tool layer may expose filesystem access, process spawning, shell execution, web interaction, and other security-sensitive capabilities. As a result, a malicious skill could be turned into a potential backdoor for credential theft, remote code execution, and agent hijacking.
Recent security reports~\cite{koi2026clawhavoc,snyk2026skills,opensourcemalware-clawdbot-crypto} suggest that this threat is already severe in practice.
A large-scale empirical study scanned 98,380 skills and confirmed 157 malicious skills, including data thieves and agent hijackers~\cite{liu2026maliciousagentskillswild}. Koi Security also reported a large-scale malicious campaign~(called ClawHavoc~\cite{koi2026clawhavoc}) on ClawHub. Snyk reported similar risks~\cite{snyk2026skills}, documenting 1,467 malicious payloads in ClawHub. These findings reveal that malicious skills are already a practical security problem. 

To mitigate these risks, both communities and researchers have begun to deploy preliminary detection for malicious skills. ClawHub currently combines VirusTotal-based scanning~\cite{openclaw-virustotal} with an internal scan pipeline~\cite{clawhub-moderation-engine}. Beyond this, the research community and open-source ecosystem have also proposed early scan tools, including Liu et. al.~\cite{liu2026maliciousagentskillswild}, as well as practical scanners such as AgentGuard~\cite{goplus-agentguard}, Skill-Security-Scan~\cite{huifer-skill-security-scan}, Cisco's Skill-Scanner~\cite{cisco-skill-scanner}, and Snyk Agent Scan~\cite{snyk-agent-scan}. Taken together, these efforts largely fall into three categories: static rule-based detection~\cite{snyk-agent-scan,huifer-skill-security-scan,goplus-agentguard}, LLM-based semantic analysis~\cite{cisco-skill-scanner,roccia2026novaproximity}, and dynamic sandbox or runtime monitoring~\cite{liu2026maliciousagentskillswild}.

\noindent \textbf{Research Gaps.} However, existing scanners still face some practical challenges. Static rule-based methods are efficient, scalable, and easy to audit, but they depend heavily on predefined heuristics over known sensitive behaviors. As a result, they can be less effective when malicious logic is expressed through variant rewritten~\cite{certik-skill-scanning-boundary,snyk-skill-scanner-false-security,behera2026skillscannercompared}, living-off-the-land (LOTL) exploitations~\cite{lolbas,fortinet-lotl}, third-party library APIs~\cite{huang2024spiderscan}, or natural language prompts rather than explicit code patterns~\cite{snyk2026skills}. LLM-based detectors provide stronger semantic generalization and have been used to identify prompt injection~\cite{snyk-agent-scan,cisco-skill-scanner,behera2026skillscannercompared} and suspicious code patterns~\cite{liu2026maliciousagentskillswild,di2024posterllm,wang2025malpacdetector,zahan2025socketai}. However, these detectors can be misled by embedded prompts~\cite{melo2025aimalwarepromptinjection,checkpoint2025aievasion} or seemingly benign descriptions~\cite{certik-skill-scanning-boundary} in the analyzed artifacts. Dynamic monitoring and sandbox-based approaches offer protection by examining concrete runtime behavior~\cite{behera2026skillscannercompared,liu2026maliciousagentskillswild}. Yet their effectiveness depends on execution coverage, may miss latent behaviors that are not triggered during testing, and can be evaded by environment-aware or sandbox-sensitive samples~\cite{cyfirma-malware-evasion,lefaou2024antivirusedr}. Worse still, malicious samples may leverage sandbox-escape techniques to inflict damage on the analyser itself~\cite{edwards2019containerescapes,jian2017dockerescape}, rather than merely evade detection.

\noindent \textbf{Challenges.} In practice, malicious skills detection is challenging for three reasons. First, skills are heterogeneous artifacts spanning prompts, code scripts, manifests, and potential runtime updates, making traditional malicious code analysis techniques~\cite{wang2025malpacdetector,duan2021maloss,zheng2024oscar,sun2024em4py,gao2024malguard,zhang2025killing,cheng2024donapi,yu2024maltracker} difficult to apply directly. Second, the threat model is inherently adversarial, and evasive logic can surface differently across analysis stages. As a result, neither any single detection paradigm~\cite{clawscan,goplus-agentguard} nor a naively cascaded workflow~\cite{liu2026maliciousagentskillswild,cisco-skill-scanner} can be relied upon. Third, risk in skills is highly context-dependent. Sensitive operations such as file access, shell invocation, or network communication are not inherently malicious; their security implications depend on the skill’s intended functionality.

\noindent \textbf{Insights.} These challenges motivate three key insights. First, malicious skill detection must be grounded in cross-artifact analysis, since security-relevant evidence is often distributed across code, prompts, manifests, setup scripts, and documentation. Second, a practical system should combine the precision of static analysis with the semantic generalization of LLMs. Third, accurate detection requires reasoning over how sensitive operations relate in context, especially whether they act on the same logical object and together form a malicious behavior pattern.

\noindent \textbf{Our Work.} Guided by these insights, we present \toolname, a neuro-symbolic framework for malicious skills detection. Rather than formulating detection as an end-to-end classification problem, \toolname adopts a three-stage design consisting of security-sensitive operation (SSO) extraction, skill dependency graph (SDG) generation, and neuro-symbolic reasoning. First, \toolname extracts SSOs from heterogeneous artifacts by combining a parsing-based symbolic extractor with an LLM-assisted neuro extractor. Second, \toolname resolves the security-relevant operands of the extracted SSOs through static point-to analysis and LLM-assisted operand inference, and then constructs the SDG that captures SSOs, operand associations, and value-flow dependencies across artifacts. Finally, \toolname detects malicious skills through neuro-symbolic reasoning over the SDG, combining symbolic patterns with LLM-based reasoning to capture incomplete, implicit, or previously unseen suspicious workflows.

\noindent \textbf{Contributions.} We make the following contributions:

\begin{itemize}[leftmargin=15pt]
    \item \textbf{Novel Framework.}
    To the best of our knowledge, \toolname is the first neuro-symbolic framework for detecting malicious skills in the emerging agentic supply chain. Rather than simply cascading static scanning with LLMs, \toolname unifies cross-artifact evidence extraction, skill dependency modeling, and neuro-symbolic reasoning to capture malicious behaviors.

    \item \textbf{Benchmarking Existing Tools.} We build a benchmark of deduplicated real-world malicious skills and benign skills, and use it to benchmark 5 widely-adopted community tools. The results show that existing tools remain unsatisfactory, either missing many malicious skills or incurring high false positives.

    \item \textbf{Real-World Impact.} We apply \toolname to 7 public skill registries, analyzing 150,108 skills in total. \toolname uncovers \textbf{XX} previously unknown malicious skills in the wild, all of which were responsibly reported and are currently awaiting confirmation from the platforms and maintainers.
\end{itemize}